\section{Why the analogy stops at norm}
\label{sec:scope}

Costa's theorem is a statement about power.
The precoder spends no extra power on known interference because the decoder cannot tell lattice-shifted signals apart.
In a write channel, power is update norm, and our first result is exactly that statement.
Continual learning, however, asks for retention, and retention is already free for projection: an A-null write keeps every old output fixed at any norm.
The dirty-paper premise and the projection premise therefore do not compete on the primary metric.
They compete only on how much norm the new task needs, and the reviewer's predicted risk applies there.
The joint precoder's saving is borrowed from the old outputs, and it is repaid with old items whenever the write is short or later training shrinks the weights.

The modulo decoder is also the scope boundary.
With a non-periodic readout, such as regression, logits or argmax over classes, the coset \(r+M\mathbb Z^{N_B}\) becomes the single point \(r\) as \(M\to\infty\), the closest-vector step is the identity, and the A-null dirty-paper write is GPM.
Our self-test checks this limit.
A learned nonlinear encoder and decoder can recover THP-like folding for communication~\citep{ozyilkan2025dirty}, but in a network the decoder is the model itself: making its readout periodic changes the task, and the model must also be trained to be stable to whole-period moves of its outputs.

\paragraph{Limitations.}
The hidden layer is frozen and the write is a single readout update, so the comparison is between closed-form writes, not training runs.
Budgets are scaled by GPM's exact-write norm, which compresses the grid at \(S=4\) where GPM is weak; the switched write and the decay curves are post hoc.
Later training is modeled as uniform weight shrink only.
We test one lattice (\(M\mathbb Z^n\), \(M=8\)) and small task sizes (\(N_A=24\), \(N_B=16\)).
A larger network could make the joint precoder's cliff shallower, but nothing here suggests it would lift old-task accuracy above projection's.

\paragraph{Verdict.}
Dirty-paper writes, as specified, are killed: no gain over GPM on old-task accuracy, and no pre-registered gain on new-task accuracy at equal norm in \nkilled{} of \nconfigs{} configurations.
What survives is an exact accounting identity, that known old weights cost the precoder no norm, which belongs to the modulo decoder and is mostly available to projected descent on a periodic loss.
